\chapter{引言}

\section{赛题目标与总体结构}

赛题二“基于鸿蒙的透明物体实时色散渲染系统”要求参赛者在给定基线工程上完成透明物体的实时折射、色散与交互展示，并进一步体现多泡泡效果、运动模拟、完整 UI 和稳定运行能力。肥皂泡恰好同时包含透明折射、薄膜干涉、动态表面和交互反馈四类典型问题，因此非常适合作为赛题落地对象。

本技术文档将介绍本队最终决赛提交版本的核心技术路线。决赛版本在初赛基础上，将系统从"单主泡泡+装饰展示泡泡"架构扩展为完整的多泡泡交互系统：泡泡之间可以动态接触、形成共享膜、融合吸收或分离，并支持运行时增删泡泡与方向性风场。实现上，系统被组织为渲染管线、薄膜材质、几何数据、动态模拟和资源加载几个相对独立的模块，并通过 ArkTS 页面层与 Native 渲染层共同完成最终应用。

\reportfigure{1.1.jpg}{0.44}{0.28\textheight}{最终整体效果：多泡泡交互场景，包含背景折射、天空盒环境、虹彩边缘与泡泡间接触共享膜。}

\section{系统组成}

最终鸿蒙提交版本的核心源码位于：
\begin{quote}
  \path{HarmonyOS_Source_Submission/refraction/}
\end{quote}
下表列出了和实现相关的各个模块。

\begin{table}[H]
  \centering
  \caption{系统模块与职责}
  \label{tab:structure}
  \begin{tabular}{p{0.22\linewidth}p{0.30\linewidth}p{0.36\linewidth}}
    \toprule
    模块 & 主要内容 & 作用 \\
    \midrule
    页面与交互 & ArkTS 页面、XComponent、按钮与参数面板 & 承载渲染画面，处理触控输入、模式切换、状态显示和参数调节。\\
    Native 接口与管线 & NAPI、EGL、FBO、绘制顺序 & 连接 ArkTS 与 C++ 渲染层，组织多阶段折射渲染和场景合成。\\
    泡泡材质 & 折射 shader、薄膜虹彩、动态膜厚 & 实现屏幕空间折射、Kim2012/LUT/Belcour Airy 三种虹彩模式和触摸波纹。\\
    渲染封装 & 相机、模型、网格、shader 工具类 & 管理天空盒、球体、动态网格和 OpenGL ES 绘制细节。\\
    表面模拟 & DBSTT/vortex sheet 单泡泡动态 & 让主泡泡轮廓和法线随时间产生轻微形变。\\
    \midrule
    多泡泡交互 & 接触图、空间哈希、表面控制点、共享膜 & 管理泡泡间接触检测、接触圆动力学、持久表面变形、共享膜生成与弯曲、融合/分离事件。\\
    风场与动态管理 & 方向性风场、动态增删泡泡 & 提供可调节的全局风场，支持运行时按需创建和删除泡泡。\\
    资源文件 & 天空盒、薄膜 LUT、HAP 提交资源 & 提供环境反射、背景折射参照、预计算薄膜反射数据和移动端运行素材。\\
    \bottomrule
  \end{tabular}
\end{table}

\section{运行环境}

本项目初期开发在 Windows 主机上进行，验证设计可行性和效果后使用 DevEco Studio 完成基于 HarmonyOS 架构的实现。推荐环境如下：

\begin{table}[H]
  \centering
  \caption{推荐编译与运行环境}
  \begin{tabular}{p{0.28\linewidth}p{0.58\linewidth}}
    \toprule
    项目 & 要求 \\
    \midrule
    开发主机操作系统 & Windows 11。\\
    开发工具 & DevEco Studio 6.0.2。\\
    SDK 环境 & HarmonyOS SDK 6.0.2（API 22）。\\
    页面层 & ArkTS + XComponent。\\
    Native 层 & CMake + C++ + EGL + OpenGL ES 3.0 + NAPI。\\
    真机测试 & 华为 Mate 70 Pro（HarmonyOS 6.0，API 22）。\\
    模拟器 & HarmonyOS 本地模拟器。\\
    \bottomrule
  \end{tabular}
\end{table}

命令行构建可使用 Hvigor，例如：

\begin{lstlisting}[language=bash]
node "E:\huawei\DevEco Studio\tools\hvigor\bin\hvigorw.js" assembleHap --no-daemon
\end{lstlisting}

签名完成后，构建产物通常位于：
\begin{quote}
  \path{entry/build/default/outputs/default/entry-default-signed.hap}
\end{quote}
提交的文件中已经整理出 \codepath{HarmonyOS\_Source\_Submission} 下的“可执行程序”目录。
